\section{Requirements Analysis}
\subsection{Project Scope and Component Requirements}

The requirements for this project span multiple interconnected components, each
critical to the realisation of a functional embedded suspicious luggage
detection system. These components can be categorised into four primary areas:
data acquisition and preparation, model selection and optimisation, algorithmic
definition and implementation, and hardware deployment.

\subsubsection{Dataset Procurement and Preparation}
\label{sec:requirements_dataset_procurement_preparation}

A robust dataset is fundamental to training and validating an effective object
detection model. The requirements in this area include:

\begin{itemize}
	\item \textbf{Source Selection:}\label{req:source_selection} Identifying suitable publicly available
	      datasets or sourcing custom surveillance footage that reflects
	      real-world deployment scenarios. As discussed in
	      section~\ref{sec:evaluation_metrics_for_object_detection}, existing
	      abandoned object detection datasets (ABODA, I-LIDS, PETS) are limited
	      in availability and recency.
	\item \textbf{Annotation Requirements:}\label{req:annotation} Ensuring that training data includes
	      precise bounding box annotations for objects of interest (luggage,
	      people) and contextual information (ownership associations, temporal
	      sequences).
	\item \textbf{Data Diversity:}\label{req:data_diversity} The dataset must encompass various
	      environmental conditions, camera angles, lighting scenarios, and
	      crowding levels to ensure the model generalises to real deployment
	      environments.
	\item \textbf{Dataset Splitting:}\label{req:dataset_splitting} Appropriate partitioning into training,
	      validation, and test sets to enable unbiased model evaluation and
	      prevent overfitting.
\end{itemize}

In addition, a dataset specifically for testing the end-to-end pipeline with
realistic scenarios of suspicious luggage is required.

\begin{itemize}
	\item \textbf{Scenario Coverage:}\label{req:scenario_coverage} The dataset should include a variety of
	      scenarios relevant to the target application, such as luggage being
	      left unattended, moved by individuals, or in close proximity to
	      people.
	\item \textbf{Environmental Diversity:}\label{req:env_diversity} The dataset should encompass
	      different environments (indoor, outdoor), lighting conditions (daytime,
	      nighttime), and crowd densities to ensure the robustness of the system.
	\item \textbf{Number of Examples:}\label{req:num_examples} A sufficient number of examples for each scenario should
	      be included to allow for a comprehensive evaluation of the system's
	      performance.
\end{itemize}

\subsubsection{Object Detection Model Selection}

As established in the literature review, single-stage detectors, particularly
the YOLO family, are the most suitable candidates for real-time edge
deployment. Model selection must consider:

\begin{itemize}
	\item \textbf{Architecture Choice:}\label{req:arch_choice} Evaluating YOLO variants (YOLO11,
	      YOLO26) to determine which offers the optimal balance between
	      accuracy and inference speed on the target hardware (NVIDIA Jetson
	      Orin Nano).
	\item \textbf{Model Size:}\label{req:model_size} Selecting an appropriate model variant (nano,
	      small, medium) that fits within the computational and memory
	      constraints of the embedded device whilst maintaining sufficient
	      detection accuracy for small objects in surveillance footage.
	\item \textbf{Transfer Learning Strategy:}\label{req:transfer_learning} Leveraging pre-trained weights
	      from large-scale datasets (COCO) as a starting point, then fine-tuning
	      on domain-specific luggage detection data to improve performance in
	      the target application.
	\item \textbf{Performance Metrics:}\label{req:perf_metrics} Establishing baseline performance
	      expectations in terms of mAP, inference latency, and frames-per-second
	      (FPS) on the embedded device.
\end{itemize}

\subsubsection{Object Tracking Model Integration}

Maintaining consistent object identities across video frames is essential for
associating luggage with owners and detecting abandonment. This requires:

\begin{itemize}
	\item \textbf{Tracker Selection:}\label{req:tracker_selection} Choosing a real-time, lightweight tracking
	      algorithm such as SORT, DeepSORT, ByteTrack, or BoT-SORT that can
	      efficiently associate detections across frames with minimal
	      computational overhead.
	\item \textbf{Integration with Detection:}\label{req:tracker_integration} Designing the pipeline to
	      seamlessly integrate the detector and tracker, ensuring that
	      detection outputs are passed to the tracker for identity maintenance.
	\item \textbf{Robustness Requirements:}\label{req:tracker_robustness} The tracker must handle common
	      surveillance challenges including temporary occlusions, identity
	      switches, and brief object disappearances.
	\item \textbf{Performance on Edge:}\label{req:edge_performance} Ensuring the combined detection and
	      tracking pipeline meets real-time processing requirements on the
	      embedded hardware without thermal throttling during sustained
	      operation.
\end{itemize}

\subsubsection{Defining and Detecting Suspicious Luggage}

As discussed in the introduction, luggage becomes \enquote{suspicious} only in
a temporal context. Operationalising this definition requires:

\begin{itemize}
	\item \textbf{Ownership Association:}\label{req:ownership} Developing logic to associate detected
	      luggage with nearby persons, establishing ownership relationships that
	      persist across frames.
	\item \textbf{Thresholds:}\label{req:thresholds} Defining configurable parameters such as
	      the maximum distance (radius) between an owner and their luggage, and
	      the minimum duration of separation before luggage is flagged as
	      abandoned.
	\item \textbf{State Machine Design:}\label{req:state_machine} Implementing a state machine that
	      tracks luggage states (attended, unattended, abandoned) and transitions
	      based on spatial and temporal conditions.
	\item \textbf{Alert Generation:}\label{req:alert_gen} Specifying the criteria and format of
	      alerts when suspicious luggage is detected, ensuring operators receive
	      actionable information.
\end{itemize}

\subsubsection{Embedded Hardware Implementation and Optimisation}

Deploying the complete system on the NVIDIA Jetson Orin Nano requires
addressing hardware-specific constraints:

\begin{itemize}
	\item \textbf{Hardware Characterisation:}\label{req:hw_char} Profiling the Jetson Orin Nano to
	      understand its computational capabilities, memory constraints, and
	      thermal characteristics.
	\item \textbf{Model Optimisation:}\label{req:model_opt} Applying techniques such as TensorRT
	      quantisation (FP16) and kernel auto-tuning to maximise inference speed
	      whilst preserving accuracy.
	\item \textbf{Memory Management:}\label{req:memory} Ensuring the complete system (detector,
	      tracker, abandonment logic) operates within the device's LPDDR5 memory
	      constraints (8 GB).
	\item \textbf{Power Efficiency:}\label{req:power} Optimising for sustained operation under
	      typical power budgets, avoiding thermal throttling and performance
	      degradation during extended operation.
	\item \textbf{Software Framework Selection:}\label{req:framework} Selecting appropriate deep
	      learning frameworks and deployment tools (PyTorch with ONNX export,
	      TensorRT, and NVIDIA DeepStream SDK) that balance ease of development
	      with runtime efficiency.
\end{itemize}

\subsection{Requirements Traceability}
\label{sec:requirements_traceability}

Table~\ref{tab:traceability} (appendix) maps each requirement to the section where it is
addressed in the implementation and evaluation, providing traceability between
the project's goals and their realisation.

